\documentclass[a4paper,twoside,12pt]{article}
\usepackage[top=1.5cm,left=1.5cm,right=1.5cm,bottom=2cm]{geometry}

\usepackage[utf8]{inputenc}
\usepackage[T1]{fontenc}
\usepackage{lmodern}
\usepackage[brazilian]{babel}

\usepackage{graphicx}
\usepackage{subfig}
\usepackage{psfrag}
\usepackage{cite}
\usepackage[cmex10]{amsmath}
\usepackage{amssymb}
\usepackage{upgreek}
\usepackage{microtype}
\usepackage{titling}
\usepackage{courier}
\usepackage{url}
\usepackage{hyperref}
\hypersetup{
	colorlinks=true,
	linkcolor=blue,
	urlcolor=blue
}

% Use o pacote 'listings' para adicionar código e o pacote 'color' para gerar os destaques
\usepackage{listings}
\usepackage{color}
\definecolor{grey}{rgb}{0.6,0.6,0.6}
\definecolor{codegreen}{rgb}{0,0.6,0}
\definecolor{codegray}{rgb}{0.5,0.5,0.5}
\definecolor{codepurple}{rgb}{0.58,0,0.82}
\definecolor{backcolour}{rgb}{0.95,0.95,0.92}
\lstset{language=Python,				% definir para destaque de linguagem Python
	backgroundcolor=\color{backcolour},
	commentstyle=\color{codegreen},
	keywordstyle=\color{magenta},
	numberstyle=\tiny\color{codegray},
	stringstyle=\color{codepurple},
	basicstyle=\ttfamily\small,		% definir tamanho das fontes
	breaklines=true,				% definir quebra de linha automática
	linewidth=\textwidth,			% definir tamanho da caixa de código
	showstringspaces=false}			% mostrar espaços como sublinhados apenas dentro de strings


\title{\vspace{-2cm} Processamento Digital de Sinais\\ Tutorial 11 - Periodograma e Método de Welch}
\author{\url{https://github.com/fchirono/Aulas_PDS_Acustica}}

\begin{document}
	\date{}
	\maketitle
	
	\section*{Objetivo de Aprendizagem}
	Ao final desta sessão, você será capaz de:
	\begin{itemize}
		\item Estimar a Densidade Espectral de Potência (DEP) de uma sequência em tempo discreto usando o método do periodograma, e reconhecer suas limitações;
		\item Estimar o Espectro de Potência (EP) e a Densidade Espectral de Potência (DEP) de uma sequência em tempo discreto usando o método de Welch;
		\item Compreender a diferença entre as normalizações de Espectro de Potência e Densidade Espectral de Potência, e escolher a normalização apropriada de acordo com a natureza determinística ou aleatória do sinal analisado;
		\item Avaliar o efeito do comprimento da DFT ($N_{DFT}$) e do número de quadros médios sobre a resolução em frequência e a variância das estimativas espectrais.
	\end{itemize}
	
	\begingroup
	\let\clearpage\relax
	\tableofcontents
	\endgroup
	
	% *-*-*-*-*-*-*-*-*-*-*-*-*-*-*-*-*-*-*-*-*-*-*-*-*-*-*-*-*-*-*-*-*-*-*-*-*-*-*-*-*-*-*-*-
	% *-*-*-*-*-*-*-*-*-*-*-*-*-*-*-*-*-*-*-*-*-*-*-*-*-*-*-*-*-*-*-*-*-*-*-*-*-*-*-*-*-*-*-*-
	\section{Revisão}
	\subsection{Periodograma}
	
	Seja $x[n]$, $n = 0, \ldots, N-1$ uma sequência em tempo discreto. O periodograma é a estimativa espectral obtida a partir da DFT normalizada de um único segmento (ou ``quadro'') de comprimento $N$ desta sequência, aplicando-se opcionalmente uma janela $w[n]$:
	
	\begin{equation}
		X_w(\Omega) = \sum_{n=0}^{N-1} w[n] x[n] e^{-j \Omega n}.
	\end{equation}
	
	Diferentemente do estimador de Welch (ver Tutorial 10), o periodograma \textbf{não} reduz a variância da estimativa conforme a quantidade de dados analisados aumenta: ao aumentar o comprimento $N$ do sinal analisado, a resolução em frequência $\Delta f = f_s / N$ melhora (i.e. diminui), mas a variância da estimativa permanece constante. Como consequência, o periodograma não converge para a Densidade Espectral de Potência teórica do sinal conforme mais dados são adicionados.
	
	% *-*-*-*-*-*-*-*-*-*-*-*-*-*-*-*-*-*-*-*-*-*-*-*-*-*-*-*-*-*-*-*-*-*-*-*-*-*-*-*-*-*-*-*-
	\subsection{Normalizações do Espectro de Potência e da Densidade Espectral de Potência}
	
	Para uma sequência $x[n]$ de comprimento $N$, janelada por $w[n]$, e com fatores de normalização de janela $S_1 = \sum_{n=0}^{N-1} w[n]$ e $S_2 = \sum_{n=0}^{N-1} w^2[n]$, as duas normalizações possíveis do estimador (periodograma ou Welch) são:
	
	\begin{itemize}
		\item \textbf{Espectro de Potência}: resulta em estimativas consistentes para a amplitude (ao quadrado) de sinais deterministicos (ex: senoides), independentemente do comprimento $N_{DFT}$ usado:
		\begin{equation}
			EP(\Omega) = \frac{2}{S_1^2} \left| X_w(\Omega) \right|^2;
		\end{equation}
		
		\item \textbf{Densidade Espectral de Potência}: resulta em estimativas consistentes para a densidade espectral de potência de sinais não-deterministicos (ex: ruído), independentemente do comprimento $N_{DFT}$ usado:
		\begin{equation}
			DEP(\Omega) = \frac{2}{f_s S_2} \left| X_w(\Omega) \right|^2.
		\end{equation}
	\end{itemize}
	
	\noindent O fator ``2'' em ambas as expressões é necessário para compensar a energia/potência contida na banda unilateral não calculada (i.e. as frequências negativas), de modo que a energia/potência total do sinal seja preservada.
	
	A potência de um sinal determinístico pode ser recuperada somando-se todos os valores do Espectro de Potência. Já a potência de um sinal aleatório dentro de uma dada faixa de frequências pode ser recuperada integrando-se (i.e. somando e multiplicando pela resolução em frequência $\Delta f$) a Densidade Espectral de Potência dentro desta faixa.
	
	% *-*-*-*-*-*-*-*-*-*-*-*-*-*-*-*-*-*-*-*-*-*-*-*-*-*-*-*-*-*-*-*-*-*-*-*-*-*-*-*-*-*-*-*-
	% *-*-*-*-*-*-*-*-*-*-*-*-*-*-*-*-*-*-*-*-*-*-*-*-*-*-*-*-*-*-*-*-*-*-*-*-*-*-*-*-*-*-*-*-
	\section{Tarefas}
	\begin{enumerate}
		\item Gere um sinal do tipo ruído branco {\tt x\_ruidobranco} com $T = 2$ s de duração, amostrado a $f_s = 48000$ Hz, com distribuição normal (i.e. gaussiana) de média zero e desvio padrão unitário. Filtre este sinal usando um filtro passa-faixas Butterworth de 4\textsuperscript{a} ordem, com frequências de corte $f_{c,baixa} = 50$ Hz e $f_{c,alta} = 2000$ Hz, obtendo o sinal {\tt x\_passafaixa}. Compare graficamente os dois sinais no domínio do tempo.
		
		\item Escreva uma função Python {\tt periodograma(x, fs, nome\_janela, normalizacao)} que calcule o periodograma unilateral de um sinal {\tt x}, com opção de escolha da janela (use {\tt scipy.signal.get\_window()}) e da normalização (``espectro'' ou ``densidade''). A função deve retornar o vetor de frequências {\tt f} e a estimativa espectral correspondente.
		
		\item Usando a função {\tt periodograma}, calcule a Densidade Espectral de Potência do sinal {\tt x\_passafaixa} considerando diferentes durações do sinal (por exemplo, $N_{per}$, $2 N_{per}$, $5 N_{per}$ e $10 N_{per}$ amostras, com $N_{per} = 512$). Plote as curvas obtidas em um mesmo gráfico (eixo de frequência logarítmico) e comente o efeito do aumento da duração do sinal sobre a resolução em frequência e sobre a variância da estimativa.
		
		\item Escreva uma função Python {\tt welch(x, Ndft, N\_sobreposicao, nome\_janela, fs, normalizacao)} que calcule a estimativa espectral de {\tt x} pelo método de Welch, com sobreposição e janelamento configuráveis, e com opção de escolha da normalização (``espectro'' ou ``densidade''). A função deve retornar o vetor de frequências {\tt f} e a estimativa espectral correspondente.
		
		\item Usando a função {\tt welch}, calcule a Densidade Espectral de Potência do sinal {\tt x\_passafaixa} mantendo $N_{DFT} = 512$ e $N_{sobreposicao} = 256$ fixos, mas variando a duração total do sinal analisado (por exemplo, 1000, 2000, 5000 e 10000 amostras). Plote as curvas obtidas em um mesmo gráfico e compare o resultado com o obtido na Tarefa 3 usando o periodograma. O que muda em relação à resolução em frequência e à variância da estimativa?
		
		\item Gere um sinal senoidal {\tt x\_senoidal} de frequência $f_0 = 1234$ Hz e amplitude RMS $V_{rms} = 2$ V, com a mesma duração e frequência de amostragem do sinal {\tt x\_passafaixa}. Some os dois sinais para obter {\tt x\_soma = x\_senoidal + x\_passafaixa}.
		
		\item Usando a função {\tt welch} com janela de Hann e normalização ``densidade'', calcule a DEP do sinal {\tt x\_soma} para diferentes comprimentos de DFT (por exemplo, $N_{DFT}$, $2 N_{DFT}$, $5 N_{DFT}$ e $10 N_{DFT}$ amostras, mantendo a mesma proporção de sobreposição). Plote as curvas obtidas (eixos log-log) e comente o comportamento do pico correspondente à senoide e do patamar correspondente ao ruído passa-faixa, conforme $N_{DFT}$ varia.
		
		\item Repita a Tarefa 7 usando agora a normalização ``espectro'' (Espectro de Potência) ao invés de ``densidade''. Compare o valor do pico obtido com $V_{rms}^2$ do sinal senoidal, para os diferentes valores de $N_{DFT}$ utilizados. O resultado é consistente independentemente do valor de $N_{DFT}$? Compare também o comportamento do patamar de ruído nesta normalização com o observado na Tarefa 7.
		
		\item Com base nos resultados obtidos, explique por que as curvas de DEP/EP calculadas pelo método de Welch (Tarefas 7 e 8) apresentam diferentes graus de variância (i.e. de ``rugosidade'') para diferentes valores de $N_{DFT}$, dado que a duração total do sinal analisado é mantida constante.
	\end{enumerate}
	
\end{document}